\chapter{认知复杂度治理的五种基本动作}
\label{chap:cognitive-complexity-governance-five-actions}

\begin{center}
\fbox{
\parbox{0.92\textwidth}{
\textbf{本章总说明}

本章讨论 AgentOS 在工作记忆有限、认知资源有限以及现实输入持续开放的条件下，维持认知系统可运行性的五种基本治理动作：\textbf{截流、压缩、时序展开、迁移与重构}。这五种动作最初是在“认知复杂度治理”框架中提出的，用于回答一个直接的工程问题：当一个智能体面对的问题结构超过当前工作记忆所能同时稳定承载的范围时，它除了继续增加计算量之外，还能够做什么？

随着本书理论从“复杂度迁移”发展到 CSO 与 Agent-CSO 的认知结构本体，本章对五种动作作进一步严格化：真正被操作的并不是某种作为实体存在的“复杂度”，而是进入智能体内部的 Agent-CSO 集合、它们之间的关系结构、激活状态、处理顺序、功能位置以及承载边界。“复杂度”在本章中仅表示这些结构给当前认知系统造成的运行负载。因此，五种治理动作可以统一理解为：\textbf{改变认知结构进入系统的数量、结构表达的密度、同时激活的时间关系、结构所在的承载位置以及问题本身的表示方式，使有限工作记忆始终只承担当前真正需要全局协调的认知结构。}

本章仍然属于固定功能拓扑 AgentOS 的原理论。这里所谓“重构”，首先指问题表示、任务结构、环境关系和局部认知组织方式的改变，并不预设智能体已经拥有任意修改自身核心功能拓扑的能力。更深层的拓扑发育问题将在本书更高层的结构动力学中讨论。
}}
\end{center}

AgentOS 的工作记忆有限容量定理给出了本章的直接出发点。设时刻 $t$ 进入显式认知面的 Agent-CSO 构成图

\[
\mathcal G_h(t)
=
\left(
V_h(t),E_h(t)
\right),
\]

其中 $V_h$ 表示当前显式激活的认知结构对象，$E_h$ 表示这些对象之间当前必须同时维持的语义、因果、约束、时间、目标或动作关系。我们不把工作记忆负载简单定义为对象数量，而定义一个更一般的结构负载泛函

\[
\mathcal C_h(t)
=
\Phi
\left(
|V_h|,
|E_h|,
D_h,
U_h,
K_h,
R_h
\right),
\]

其中 $D_h$ 表示依赖深度，$U_h$ 表示不确定性，$K_h$ 表示结构耦合强度，$R_h$ 表示当前认知结构之间的冲突与竞争程度。$\Phi$ 并不是宇宙中的某种物理“复杂度能量”，而是用于预测系统是否接近认知容量边界的工程测度。

工作记忆稳定运行要求

\[
\boxed{
\mathcal C_h(t)
\leq
\mathcal C_{\max}(h)
}
\]

至少在绝大多数正常运行时间内成立。

因此，AgentOS 的核心问题并不是如何无限提高 $\mathcal C_{\max}$，而是如何主动改变进入 $\mathcal G_h$ 的结构，使系统能够在有限容量中持续完成开放世界任务。我们将在本章看到，截流、压缩、时序展开、迁移和重构分别作用于五个不同的位置：

\[
\boxed{
\text{进入什么}
\rightarrow
\text{如何表示}
\rightarrow
\text{何时处理}
\rightarrow
\text{在哪里处理}
\rightarrow
\text{以什么问题结构处理}
}
\]

它们共同形成 AgentOS 最基础的认知治理算子集合

\[
\boxed{
\mathfrak G
=
\{
\mathcal A,
\mathcal K,
\mathcal T,
\mathcal M,
\mathcal R
\}
}
\]

其中

\[
\mathcal A
=
\text{Admission},
\qquad
\mathcal K
=
\text{Compression},
\]

\[
\mathcal T
=
\text{Temporalization},
\qquad
\mathcal M
=
\text{Migration},
\qquad
\mathcal R
=
\text{Reorganization}.
\]

我们下面逐一讨论这五种动作，并特别关注它们之间并非并列的工具集合，而是构成了一条由浅入深的认知治理层级。

\section{截流：阻止不必要的认知结构进入工作记忆}

截流是五种认知复杂度治理动作中代价最低、发生位置最靠前，同时也是最容易被低估的一种。我们首先需要建立一个基本认识：面对开放世界的 Agent 不可能，也没有必要，把所有能够观测到的信息都转化成当前需要显式处理的 Agent-CSO。世界中的潜在结构数量远大于工作记忆能够稳定承载的结构数量，因此一个成熟智能系统的第一项能力不是“能够理解越来越多东西”，而是能够稳定判断“什么此刻不需要被理解”。

设外部世界和内部系统在时刻 $t$ 产生候选认知结构集合

\[
\mathcal X_t
=
\{
C_1,C_2,\ldots,C_n
\},
\]

其中每个 $C_i$ 在进入 Agent 以后可能形成一个候选 Agent-CSO

\[
\pi_A(C_i).
\]

如果我们简单执行

\[
\mathcal X_t
\rightarrow
h(t),
\]

那么工作记忆负载将随外界刺激规模直接增长。这种结构本身无法扩展。现实世界中的传感器、网络、其他 Agent、工具、文件系统以及 Agent 自身内部状态都可能持续产生数据，而这些数据的生成速度原则上可以远高于显式认知所能够处理的速度。因此，必须首先存在一个准入算子

\[
\mathcal A_t:
\mathcal X_t
\rightarrow
\mathcal X_t^{+},
\]

其中

\[
\mathcal X_t^{+}
\subseteq
\mathcal X_t
\]

代表允许进入当前显式处理面的结构集合。

我们可以进一步定义每个候选结构的准入函数

\[
a_i(t)
=
A
\left(
Rel_i,
Goal_i,
Risk_i,
Novel_i,
Residual_i,
Uncertainty_i,
Cost_i
\right),
\]

分别描述该结构与当前目标的相关度、潜在价值、风险程度、新颖性、异常残差、不确定性以及进入工作记忆所造成的处理成本。只有当

\[
a_i(t)
\geq
\theta_A
\]

时，结构才需要以完整或近完整形式进入当前全局工作空间。

因此，截流的本质不是一个简单的“过滤器”，而是一个认知准入控制问题：

\[
\boxed{
AdmissionControl
=
DecideWhatDeservesGlobalCognition
}
\]

这与传统计算机系统中的 admission control 有一定形式相似性：服务器不会因为存在无限请求就同时运行全部请求，实时系统也不会让所有事件获得相同调度优先级。但 AgentOS 的截流比传统任务准入更复杂，因为它处理的不是预先定义好的进程，而是具有语义、价值和未来影响的认知结构。

读者尤其需要注意，截流并不等同于“忽略”。至少可以存在四种结果。第一类结构被完整提升进入 $h$；第二类结构只形成低分辨率摘要；第三类结构保持在局部功能或外部存储中，只保留一个可恢复索引；第四类结构在当前情况下被真正丢弃。因此更准确地写，应有

\[
\mathcal A(C_i)
\in
\{
Full,
Summary,
Indexed,
Discarded
\}.
\]

成熟智能真正重要的能力，是让绝大多数正常世界动力学停留在局部闭环中，而不是不断把它们升级成显式思想。一个熟练驾驶者并不会持续把发动机每一次振动、每一个视觉像素以及每一个方向盘微调送入完整显式推理；一个成熟数字 Agent 同样不应把每次文件 IO、API 返回字段或窗口像素都变成工作记忆中的显式认知对象。只有当某些结构产生了持续残差、风险、新颖性、目标关联或者跨局部系统的冲突时，它们才应被提升。

因此我们可以得到一个非常重要的成熟智能原则：

\[
\boxed{
MostDynamicsShouldRemainLocal
}
\]

而：

\[
\boxed{
RelevantDeviationBecomesExplicitCognition.
}
\]

如果定义局部动力学产生的残差为 $\varepsilon_i(t)$，那么一种典型事件触发准入形式可以写为

\[
\Gamma_i(t)
=
w_1|\varepsilon_i|
+
w_2 Risk_i
+
w_3 Novel_i
+
w_4 Goal_i
+
w_5 Persistence_i.
\]

当

\[
\Gamma_i(t)<\theta_A
\]

时，

\[
C_i
\rightarrow
LocalClosure,
\]

而当

\[
\Gamma_i(t)\geq\theta_A
\]

时，

\[
C_i
\rightarrow
h.
\]

这种结构使工作记忆不再是所有信息流必经的“中央总线”，而成为一个稀疏的跨尺度显式协调面。我们由此也能够理解为什么工作记忆有限并不天然意味着智能不足。恰恰相反，如果系统能够通过可靠的截流让绝大多数已掌握、低风险、局部可闭合的过程停留在低层，那么有限工作记忆反而迫使系统形成更优良的层级结构。

从外部理论看，这一机制可以与选择性注意、event-triggered control、实时系统的 admission control、信息瓶颈思想以及预测加工中的 prediction error gating 进行比较。但本章不把其中任何一个理论直接等同于 AgentOS。我们真正关心的是一个更一般的问题：在有限认知容量下，一个持续 Agent 必须建立一种结构性机制，使“是否进入显式认知”本身成为可调度、可学习并受目标与现实反馈约束的决策变量。

因此，五种治理动作中的第一条原则是：

\[
\boxed{
\text{能不进入工作记忆的问题，首先不要让它进入工作记忆。}
}
\]

截流并不是降低智能，而是防止有限显式认知被那些本可以在其他尺度闭合的过程占满。只有在这一前提成立以后，后面的压缩、时序展开、迁移和重构才真正有意义。

\section{压缩：减少被保留认知结构的显式自由度}

即使经过截流，仍然会存在一批必须进入当前认知面的结构。此时第二个问题不是继续删除信息，而是：对于确实需要保留的认知结构，我们是否必须以原始粒度同时承载其全部细节？

答案通常是否定的。这就产生第二种治理动作：压缩。

我们首先把压缩与截流严格区分。截流改变的是

\[
\boxed{
WhetherAStructureEnters
}
\]

而压缩改变的是

\[
\boxed{
HowTheStructureIsRepresentedAfterItEnters.
}
\]

假设原始 Agent-CSO 子图为

\[
\mathcal G
=
(V,E),
\]

压缩算子

\[
\mathcal K:
\mathcal G
\rightarrow
\widetilde{\mathcal G}
\]

生成一个较低自由度的表示，使得

\[
|\widetilde V|
<
|V|
\]

或

\[
|\widetilde E|
<
|E|,
\]

但这种“更小”本身并不是压缩成功的充分条件。真正要求的是：对当前决策有意义的关系应尽可能被保存。

因此，我们可以把认知压缩写成一个约束优化问题：

\[
\boxed{
\min_{\widetilde{\mathcal G}}
\;
\mathcal C(\widetilde{\mathcal G})
}
\]

subject to

\[
\mathcal L_{task}
\left(
\mathcal G,
\widetilde{\mathcal G}
\right)
\leq
\epsilon.
\]

其中 $\mathcal C$ 表示工作记忆中维持这一结构所需要的有效负载，而 $\mathcal L_{task}$ 表示压缩前后对于当前目标、预测和行动而言损失了多少关键结构。

这种形式说明，认知压缩从来不是简单“摘要得越短越好”。一个只有一句话的摘要如果丢失了决定行动所需要的关键条件，其长度虽然极小，却是失败的压缩。相反，一个相对较长的结构如果保留了当前所有关键约束，反而可能具有更好的认知效率。

从信息论语言看，我们可以借用“充分统计量”的思想。设当前任务相关变量为 $Y$，原始认知结构为 $X$，压缩结构为 $Z$：

\[
X
\xrightarrow{\mathcal K}
Z.
\]

理想情况下，我们希望：

\[
I(Z;Y)
\approx
I(X;Y)
\]

同时：

\[
H(Z)
\ll
H(X).
\]

这可以被理解成：压缩结构对于当前行动仍保留接近完整的任务相关信息，但自身自由度显著下降。这里使用 mutual information 只是给出一种数学解释，并不意味着真实 AgentOS 必须显式计算精确的信息熵。

在实际认知过程中，压缩可能至少表现为五种不同方式。

第一种是对象聚合。许多低层对象可以形成一个更高层类别。例如：

\[
\{
car_1,
car_2,
car_3,
\ldots
\}
\rightarrow
TrafficFlow.
\]

第二种是关系折叠。大量重复关系可以被一个结构规则替代，例如：

\[
A_1\to B_1,\quad
A_2\to B_2,\quad
\ldots
\]

被压缩成：

\[
A_i
\xrightarrow{Rule}
B_i.
\]

第三种是情景摘要。一个完整事件轨迹

\[
E=
(x_1,x_2,\ldots,x_n)
\]

可以在当前任务中只保留关键转折点

\[
\widetilde E=
(x_{k_1},x_{k_2},\ldots,x_{k_m}),
\qquad
m\ll n.
\]

第四种是子目标折叠。已经解决的一系列步骤不必长期作为工作记忆的活动节点存在，而可以重新折叠成“该子问题已满足”这一更高层状态。

第五种是结构命名。当一组关系获得稳定名称之后，一个符号本身便可以成为对复杂关系网络的索引。这也是语言、数学符号和专业概念为什么能够显著扩大人类有效认知能力的重要原因之一。

因此，压缩不是简单减少数据量，而是在改变 Agent-CSO 的表示坐标：

\[
AgentCSO
=
(C,R,F,\tau,B,S)
\]

中的

\[
R:
R_{high-dimensional}
\rightarrow
R_{compressed}.
\]

本体结构 $C$ 可以保持相对连续，而具体显式 representation $R$ 被重新组织。

这里还需要讨论压缩的一个核心风险：不可逆信息损失。

如果压缩函数 $\mathcal K$ 过于激进，则可能产生：

\[
\mathcal K(C_1)
=
\mathcal K(C_2)
\]

即两个实际上影响未来行动不同的结构被错误合并。结果是 Agent 在当前状态下认为二者等价，而现实随后产生不同反馈。我们可以将这种风险定义成 compression-induced ambiguity：

\[
A_{\mathcal K}
=
P
\left(
Y(C_1)\neq Y(C_2)
\mid
\mathcal K(C_1)=\mathcal K(C_2)
\right).
\]

因此压缩强度必须与风险水平相关：

\[
Risk\uparrow
\Rightarrow
CompressionAggressiveness\downarrow.
\]

例如，在轻松的日常任务中，Agent 可以用高度摘要的状态；在高安全场景中，却必须保留更精细的结构。

这说明认知压缩本身就是一个控制变量，而不是静态编码格式。一个成熟系统需要根据目标、风险、不确定性、未来可恢复性不断调整 representation resolution。可以定义压缩率

\[
\rho_K
=
\frac{\mathcal C(\widetilde{\mathcal G})}
{\mathcal C(\mathcal G)},
\qquad
0<\rho_K\leq1,
\]

但真正优化的并不是单独最小化 $\rho_K$，而是平衡：

\[
J_K
=
\lambda_C\mathcal C
+
\lambda_E\mathcal L_{task}
+
\lambda_R Risk
+
\lambda_U Uncertainty.
\]

从外部理论看，这一部分可以与信息瓶颈、最小描述长度、充分统计量、representation learning、chunking、schema formation 以及程序编译中的中间表示优化进行比较。但 AgentOS 关注的是它们在持续认知运行时中的共同作用：系统必须不断把已经进入显式工作面的结构重新折叠成更适合当前目标的形态。

因此第二条治理原则可以表达为：

\[
\boxed{
\text{必须进入工作记忆的结构，也不应以比当前任务需要更高的自由度存在。}
}
\]

截流解决“进不进来”，压缩解决“进来以后占多大”。两者结合以后，AgentOS 才能够真正开始控制当前认知面的结构密度。

\section{时序展开：把同时结构负载转换为时间结构}

即使一个问题中的所有重要结构都必须保留，而且进一步压缩会造成不可接受的信息损失，工作记忆仍然可能过载。这时我们不能继续删除，也不能简单摘要。第三种治理动作由此出现：把原本必须同时维持的结构依赖，重新安排成一个可以逐步处理的时间过程。

我们把这种机制称为时序展开。

它解决的核心问题可以概括为：

\[
\boxed{
SimultaneousComplexity
\rightarrow
TemporalDepth.
}
\]

考虑一个认知任务图

\[
\mathcal G_T
=
(V,E),
\]

如果为了完成任务，系统在某一个时刻必须同时维持大量节点，则峰值工作记忆需求近似受到图的有效并行宽度约束。记为

\[
W(\mathcal G_T)
\]

表示某种最小必要同时活动宽度，而：

\[
D(\mathcal G_T)
\]

表示完成该任务需要经历的时间深度。

一个重要的工程思想是，很多问题存在：

\[
W
\leftrightarrow
D
\]

的交换关系。

通过将同时并行处理转换成顺序处理，可以使

\[
W'\ll W
\]

但通常代价是

\[
D'>D.
\]

因此我们可以把时序展开写成：

\[
\mathcal T:
\mathcal G_T
\rightarrow
\mathcal G_T',
\]

使：

\[
W(\mathcal G_T')
<
W(\mathcal G_T),
\]

同时允许：

\[
D(\mathcal G_T')
>
D(\mathcal G_T).
\]

这就是一种典型的“空间复杂度向时间复杂度迁移”，但按照本书后期的本体语言，更准确地说，它是：

\[
\boxed{
AgentCSOActivationStructure
:
Parallel
\rightarrow
Sequential.
}
\]

也就是说，我们改变的是认知结构的同时激活关系，而不是让某种复杂度物质从空间流向时间。

我们可以用一个简单例子理解。假设 Agent 要完成一个复杂建设任务，其中包括结构设计、材料需求、物流安排、施工顺序、风险控制和预算约束。如果系统试图在每一次决策中同时完整激活所有关系，则：

\[
\mathcal C_h
\gg
\mathcal C_{\max}.
\]

一种更合理的方法是将任务重排成：

\[
Goal
\rightarrow
Phase_1
\rightarrow
Checkpoint_1
\rightarrow
Phase_2
\rightarrow
Checkpoint_2
\rightarrow
\cdots.
\]

每一阶段只保留当前局部问题，同时把已经完成的阶段压缩成 checkpoint，把暂时未激活的未来结构留在 E、$\Theta$ 或外部计划表示中。这样，全任务仍然保持完整，但任意时刻真正进入 $h$ 的只是一个有限窗口。

因此，时序展开实际上引入了一个活动窗口：

\[
\mathcal W_t
\subset
\mathcal G_T,
\]

使：

\[
\mathcal C(\mathcal W_t)
\leq
\mathcal C_{\max}
\]

对所有阶段尽量成立，而：

\[
\bigcup_t
\mathcal W_t
\approx
\mathcal G_T.
\]

这与计算机科学中的分页、流式处理、迭代算法、动态规划和 coroutine 有一定形式对应。一个巨大数据集无法同时放入内存时，我们分批加载；一个巨大问题无法同时保存在认知空间时，Agent 同样可以进行结构分页。

但认知时序展开比普通分页多出一个重要问题：阶段之间不是独立的。之前的认知结果会改变后续问题，因此必须维护足够的跨阶段不变量。

设阶段 $k$ 的状态为

\[
z_k,
\]

阶段转换：

\[
z_{k+1}
=
F_k(z_k,u_k),
\]

则每一阶段完成后必须产生一个足够支持后续阶段的结构摘要：

\[
s_k
=
\Sigma(z_k).
\]

下一阶段并不需要完整重新加载 $z_k$，而只需要：

\[
s_k.
\]

这使时序展开与上一节的压缩自然耦合：

\[
\boxed{
Temporalization
+
CheckpointCompression.
}
\]

如果缺少这种机制，时序展开只会把原来的并行过载变成频繁重新加载和上下文丢失。

因此，AgentOS 中的时序展开必须同时解决三个问题。

第一，如何分段。不同任务需要找到近似自然闭合的子结构，使跨段依赖尽量少。

第二，如何保存阶段边界。每个阶段结束后需要生成 checkpoint、摘要、未决问题与必要约束。

第三，什么时候重新展开。如果后续阶段发现旧假设错误，系统必须能够回滚并重新激活之前被折叠的结构。

因此可以定义阶段 $k$ 的认知状态：

\[
Q_k
=
(
Goal_k,
ActiveCSO_k,
Constraints_k,
Checkpoint_{k-1},
Pending_k
).
\]

只有 $Q_k$ 进入当前显式工作面，而完整任务结构保存在更慢或外部载体中。

时序展开同样不是没有代价。它会增加：

Latency；

Checkpoint Cost；

Context Switching；

Error Propagation；

Long-Horizon Drift。

如果一个任务原本需要毫秒级响应，我们就不能通过把它展开成数分钟来解决工作记忆不足。因此可以写出：

\[
J_T
=
\lambda_W W
+
\lambda_D D
+
\lambda_S SwitchCost
+
\lambda_E ErrorPropagation.
\]

最佳时序展开不是无限序列化，而是在认知峰值和完成时间之间寻找可接受平衡。

这也说明，在多时间尺度智能系统中，工作记忆容量限制会自然催生“过程”。过程并非外部项目管理强加给智能体的一套行政规则，而是有限认知系统处理超出瞬时承载能力问题时的一种基础动力学机制。很多复杂任务之所以需要计划、阶段、里程碑、检查点和子目标，根本原因之一就在于：

\[
\boxed{
WholeProblemCannotExistExplicitlyAtOnce.
}
\]

因此第三条治理原则可以写成：

\[
\boxed{
\text{当重要结构不能继续删除和压缩时，不要要求它们同时存在；让它们在时间中逐步展开。}
}
\]

至此我们已经形成前三层治理：

\[
\boxed{
截流
\rightarrow
压缩
\rightarrow
时序展开
}
\]

它们仍然主要在当前 Agent 内部解决问题。然而，当一个问题即使被分段以后，仍然持续占据大量认知资源，我们就必须开始考虑第四种更深的动作：改变认知结构的承载位置。

\section{迁移：改变认知结构的功能位置与承载边界}

迁移是五种治理动作中第一次真正改变“认知结构在哪里存在”的操作。

截流决定结构是否进入当前工作面，压缩改变结构的表示密度，时序展开改变结构在时间上的激活方式，而迁移进一步问：

\[
\boxed{
\text{为什么这个结构必须由当前认知核心亲自承载？}
}
\]

这是 AgentOS 从传统“所有能力都在模型内部”走向广义认知系统的重要一步。

在早期理论中，我们曾把这一过程称为“复杂度卸载”或者“复杂度迁移”。按照本书现在采用的 CSO 本体，更准确地说，迁移发生在 Agent-CSO 的承载坐标上。

再次写出：

\[
AgentCSO_A(C,t)
=
(C,R,F,\tau,B,S).
\]

迁移可以改变：

\[
F:
F_i\rightarrow F_j,
\]

即功能迁移；

也可以改变：

\[
B:
Internal\rightarrow External,
\]

即边界迁移；

还可能改变：

\[
\tau:
\tau_{fast}
\rightarrow
\tau_{slow},
\]

即时间尺度迁移。

本章最关心的是承载迁移：

\[
\boxed{
\mathcal M_B:
AgentCSO(C,B_i)
\rightarrow
AgentCSO(C,B_j).
}
\]

例如，一个需要反复记住的数据表最初可能完整存在于 $h$ 中。随后系统把它写入文件：

\[
h
\rightarrow
File.
\]

未来只保留：

\[
Index(File)
\]

在内部。

一个反复进行的计算过程可以：

\[
Reasoning
\rightarrow
Script.
\]

一个稳定动作技能可以：

\[
ExplicitPlan
\rightarrow
BodySkill.
\]

一个持续约束也可以：

\[
InternalReminder
\rightarrow
EnvironmentStructure.
\]

例如，把必须记住“门不能打开”的认知要求变成物理锁，本质上就是：

\[
Constraint_{internal}
\rightarrow
Constraint_{world}.
\]

这正是认知卸载最深的含义：智能并不要求所有有效结构都持续存在于内部。只要一个结构能够通过稳定索引、调用协议和现实反馈重新参与未来行动，那么它就可以在更广义的 Agent 认知体系中继续存在。

因此可以定义广义可访问结构空间

\[
\mathcal M_A
=
\mathcal M_{internal}
\cup
\mathcal M_{tool}
\cup
\mathcal M_{file}
\cup
\mathcal M_{environment}
\cup
\mathcal M_{social}.
\]

一个结构是否属于 Agent 的有效认知资源，并不取决于它是否物理位于 Agent 内部，而取决于：

\[
\boxed{
Accessible
+
Addressable
+
Recoverable
+
Controllable.
}
\]

这一区分非常重要。

一个互联网网页虽然存在，但如果 Agent 不知道它的位置，也没有可靠检索方式，它不能简单被视为 Agent 当前可用记忆。相反，一个结构即使存在于远程数据库，只要 Agent 拥有稳定索引、认证权限、访问协议和语义映射，它就可以构成 Agent 广义外部记忆的一部分。

因此，迁移不是：

\[
InternalInformation
\rightarrow
LostInformation,
\]

而应该是：

\[
\boxed{
InternalStructure
\rightarrow
ExternalizedStructure
+
InternalIndex.
}
\]

可以将迁移代价定义为：

\[
J_M
=
C_{write}
+
C_{retrieve}
+
C_{latency}
+
C_{failure}
+
C_{dependency}
+
C_{security}.
\]

而迁移收益则主要表现为内部负载下降：

\[
\Delta C_h
=
C_h^{before}
-
C_h^{after}.
\]

只有当长期预期满足

\[
\mathbb E
[
Benefit_M
]
>
\mathbb E
[
Cost_M
]
\]

时，迁移才是合理的。

这意味着认知卸载并不是越多越好。如果所有能力都外化，Agent 可能产生巨大的工具调用依赖、延迟、安全问题和身份连续性问题。反过来，如果所有结构都坚持内部保存，又会造成持续认知膨胀。

因此真正需要优化的是：

\[
\boxed{
CarrierAllocation.
}
\]

不同类型结构应该存在于修改代价、访问速度、可靠性和持续性不同的载体。

我们可以给每种载体 $b_j$ 定义一个成本向量：

\[
Cost(b_j)
=
(
L_j,
E_j,
R_j,
M_j,
S_j
),
\]

分别表示读取延迟、能源/计算成本、失效风险、修改代价以及安全成本。对于认知结构 $C_i$，选择载体就是：

\[
b_i^\star
=
\arg\min_{b_j}
J(C_i,b_j).
\]

这进一步说明，学习不仅是“学到了什么”，还包括“学会了应该把什么放在哪里”。

这一思想解释了人类和文明中大量常见现象。人不会把所有电话号码保存在大脑里，因为联系人数据库是更好的载体；不会每次重新推导复杂公式，因为文字和数学符号保存了结果；不会每天重新规划城市交通规则，因为道路和交通信号本身把一部分规则写入现实环境。程序、工具、机构、流程甚至社会分工，都可以从不同尺度理解为认知结构寻找了更稳定、更低成本的承载体。

但本章仍然必须维持 Agent 的边界。外部结构不应拥有未经治理的直接认知控制权。任何外部返回结果仍然需要经过验证，尤其不能把：

\[
ExternalOutput
\]

自动等同于：

\[
RealityTruth.
\]

因此迁移必须与访问控制、版本控制、来源可信度和现实验证共同存在。

从外部理论看，这一机制与 cognitive offloading、extended cognition、distributed cognition、缓存层级、虚拟内存、函数调用、编译以及组织分工均存在形式对应。然而 AgentOS 更关注这些机制如何被统一纳入一个持续 Agent 的结构承载决策。

因此第四条治理原则是：

\[
\boxed{
\text{如果某个结构不值得持续占据核心认知，就应该为它寻找更合适的承载位置，而不是让工作记忆永远为已经解决的问题付费。}
}
\]

这就是从“我必须记住和计算”走向“整个世界可以帮助我记住和计算”的理论转折。

\section{重构：改变问题本身，使认知负载不再以原来的形式存在}

前四种治理动作都有一个共同前提：问题的基本表示方式仍然保持不变。

截流从问题中选择部分结构；

压缩降低结构粒度；

时序展开改变处理顺序；

迁移改变结构所在载体。

但在某些情况下，无论怎样进行这些操作，问题仍然反复造成高认知负载。这往往说明真正的问题并不是“系统资源还不够”，而是：

\[
\boxed{
TheProblemIsRepresentedInTheWrongWay.
}
\]

于是第五种，也是最深层的治理动作出现：重构。

本章中的重构主要指 representation、task decomposition、environment arrangement 和 local functional organization 的改变。它首先不是未来意义上的任意核心拓扑自修改，而是一种：

\[
\boxed{
ProblemReformulation.
}
\]

假设某个问题在表示空间 $\mathcal R_1$ 中形成认知结构

\[
\mathcal G_1
=
\Gamma(P,\mathcal R_1),
\]

其运行负载为

\[
\mathcal C(\mathcal G_1)
>
\mathcal C_{\max}.
\]

传统扩算力思想会尝试增加系统资源。但重构提出另一个可能：

寻找新的表示空间

\[
\mathcal R_2
\]

使：

\[
\mathcal G_2
=
\Gamma(P,\mathcal R_2)
\]

满足：

\[
\boxed{
\mathcal C(\mathcal G_2)
\ll
\mathcal C(\mathcal G_1)
}
\]

同时保持任务可解性。

因此：

\[
\mathcal R:
(P,\mathcal G_1)
\rightarrow
(P',\mathcal G_2).
\]

这里 $P'$ 可以在任务语义上与原问题等价，也可能通过环境改变使原问题本身变得不同。

一个经典的抽象例子是：如果 Agent 不断需要在复杂原始像素空间中推理物体关系，那么仅靠增加 context 可能永远昂贵。若建立对象中心的 SGC 式表示，则大量像素级结构可以重构为：

\[
Object
+
Relation
+
Geometry
+
Affordance.
\]

问题没有消失，但认知结构的形态发生了改变。

同样，如果一个任务需要不断重复十几个操作步骤，我们可以继续进行时序展开；但如果这些步骤已经稳定，则进一步形成一个 procedure：

\[
\{Step_1,\ldots,Step_n\}
\rightarrow
Skill.
\]

这实际上也是一种局部问题重构：原来需要面对的对象是“十几个动作”，现在只需要面对“调用某技能”。

再如，一个 Agent 每天都必须通过复杂推理判断一个环境约束是否满足。如果能够直接改变环境，使约束物理化，那么：

\[
RepeatedReasoning
\rightarrow
EnvironmentalConstraint.
\]

环境发生变化以后，原来的认知问题甚至不再存在。

因此重构可以发生至少四个层面。

第一是表示重构：

\[
Representation_1
\rightarrow
Representation_2.
\]

第二是任务重构：

\[
Problem
\rightarrow
SubproblemStructure.
\]

第三是程序重构：

\[
RepeatedReasoning
\rightarrow
ReusableProcedure.
\]

第四是环境重构：

\[
CognitiveConstraint
\rightarrow
WorldConstraint.
\]

我们可以统一定义重构收益：

\[
Gain_R
=
C_{old}
-
C_{new}
-
Cost_{transform}.
\]

但真正需要考虑的还有未来长期收益，因此可以写成：

\[
J_R
=
C_{transform}
+
\int_{t}^{t+T}
C_{new}(\tau)d\tau
-
\int_{t}^{t+T}
C_{old}(\tau)d\tau.
\]

如果：

\[
J_R<0,
\]

则说明虽然重构当下需要付出较大成本，但长期总认知成本更低。

这正是为什么重构通常应该比截流、压缩和时序展开更慢。对一个一次性的小问题重新设计整套 representation 是没有意义的。只有当某类结构反复出现、持续造成高成本、并表现出稳定规律时，重构才值得发生。

因此可以定义一个重构触发信号：

\[
\Gamma_R
=
\alpha P_{residual}
+
\beta F_{repeat}
+
\gamma C_{coord}
+
\delta L_{latency}
+
\eta C_{failure}.
\]

其中：

$P_{residual}$ 表示长期未解决残差；

$F_{repeat}$ 表示重复频率；

$C_{coord}$ 表示多结构协调成本；

$L_{latency}$ 表示由于现有表示造成的执行延迟；

$C_{failure}$ 表示现有结构造成的失败代价。

只有当：

\[
\Gamma_R>\theta_R
\]

并持续足够长时间，系统才应考虑重新组织问题表示。

这一原则非常重要，因为重构是一种高收益同时也是高风险的操作。如果系统过于频繁地重构，就会不断改变自身表示，使经验无法积累、技能无法稳定、自我和长期结构无法形成。这就是为什么 AgentOS 不能把“不断重新思考一切”视为高级智能。

真正成熟的智能恰恰相反：

\[
\boxed{
StableStructureShouldRemainStableUntilPersistentEvidenceRequiresChange.
}
\]

因此重构应具备迟滞。可以设置：

\[
\theta_{on}
>
\theta_{off}.
\]

新的结构只有在证据超过较高阈值时才建立，而一旦建立，不应因为短暂波动立即撤销。这与控制系统中的 hysteresis 以及软件工程中的 architecture refactoring 都存在形式相似性。

到这里，我们可以把五种认知治理动作统一起来。

设当前问题造成的认知负载为 $\mathcal C_0$。

首先执行截流：

\[
\mathcal C_1
=
\mathcal A(\mathcal C_0).
\]

若仍超载，则压缩：

\[
\mathcal C_2
=
\mathcal K(\mathcal C_1).
\]

若重要结构不能继续压缩，则时序展开：

\[
\mathcal C_3
=
\mathcal T(\mathcal C_2).
\]

若仍长期占据核心资源，则迁移：

\[
\mathcal C_4
=
\mathcal M(\mathcal C_3).
\]

如果问题仍然反复制造高成本，则进行重构：

\[
\mathcal C_5
=
\mathcal R(\mathcal C_4).
\]

因此形成一条由浅入深的治理阶梯：

\[
\boxed{
Admission
\rightarrow
Compression
\rightarrow
Temporalization
\rightarrow
Migration
\rightarrow
Reorganization.
}
\]

这不是一个必须机械执行的固定流水线，而是一种认知治理优先级。系统原则上应该首先采用代价更低、可逆性更强、对长期结构扰动更小的手段，只有在低层治理不足以保持系统可运行时，才逐步升级。

我们可以将这一思想称为：

\[
\boxed{
MinimumNecessaryCognitiveReorganizationPrinciple.
}
\]

即：

\[
\boxed{
\text{只进行解决当前持续结构失配所必需的最小认知重组。}
}
\]

于是本章最终给出的已经不只是五个工程技巧，而是一种完整的有限认知治理理论。

工作记忆的有限性并不意味着 Agent 必须永远受困于有限容量。真正高级的智能并不是无限扩大显式认知空间，而是持续改变什么值得进入认知、以什么粒度存在、在什么时候存在、在哪里存在，以及是否应该改变问题本身。

因此，本章可以用下面的统一表达式收束：

\[
\boxed{
\begin{aligned}
\mathcal A &: \text{改变进入认知的结构集合},\\
\mathcal K &: \text{改变结构的表示自由度},\\
\mathcal T &: \text{改变结构的同时激活关系},\\
\mathcal M &: \text{改变结构的功能与承载位置},\\
\mathcal R &: \text{改变产生认知负载的问题结构本身}.
\end{aligned}
}
\]

从 CSO 的角度看，它们实际上对应：

\[
\boxed{
Selection
\rightarrow
Representation
\rightarrow
ActivationSchedule
\rightarrow
Carrier
\rightarrow
ProblemStructure.
}
\]

而从 AgentOS 的角度看，这五种动作共同服务于同一个核心目标：

\[
\boxed{
\mathcal C_h(t)
\leq
\mathcal C_{\max}(h)
}
\]

但我们现在已经能够更加准确地理解这个公式。AgentOS 并不是简单把“复杂度降低到阈值以下”，而是在不断重新安排 Agent-CSO 与有限认知资源之间的关系。

因此，本章最核心的理论结论不是“复杂问题必须被简化”，而是：

\[
\boxed{
\textbf{
成熟智能的标志，不是它能够在工作记忆中同时容纳整个世界，而是它能够让整个世界中只有当前真正需要全局协调的结构进入工作记忆，并把其余结构稳定地安置到更适合的表示、时间和载体之中。
}
}
\]

这使工作记忆的有限性从一个缺陷转化成了智能结构成长的驱动力。正因为不能把所有东西同时放进 $h$，智能才必须形成层级、摘要、过程、技能、工具、环境结构和更加优秀的问题表示。换句话说，有限性并不只是限制智能；在相当程度上，它也是高级认知组织形式得以产生的原因。